% ===================================================================
%  اللوحة رقم ١١ — المدينة وآثارها. --- الحريق.
%
%  Traduction, faite en 2026, en arabe levantin d'aujourd'hui, sur
%  l'ido de texto/io. Regles de la colonne : voir 10-lawha-01.tex.
%  Deux parties, deux numerotations. Ouverture de la TROISIEME SERIE :
%  « السلسلة التالتة ». DEUX blocs marques « ¶2 » par ordo.py —
%  t11-c1-12-1 et t11-c2-03-1 —, et c'est le seul tableau du livret
%  qui en ait deux : une seule cle chacun, deux alineas dedans.
%
%  LE PREFET (15, 91) EST « المحافظ », ET CE N'EST PAS UNE ANALOGIE,
%  C'EST LA MEME ADMINISTRATION. La colonne gujaratie avait trouve
%  « કલેક્ટર », le Collector du district — meme charge, meme siecle,
%  mais inventee separement a Calcutta pendant qu'on l'inventait a
%  Paris. Ici le rapport est direct : la محافظة syrienne et libanaise
%  a ete decoupee sous le Mandat francais, sur le modele du
%  departement, par les gens memes dont ce livret decrit le pays. Le
%  prefet du tableau et le محافظ qui le traduit ne se ressemblent pas,
%  ILS DESCENDENT L'UN DE L'AUTRE. Et sa maison est « السراي », le
%  serail — mot turc reste sur le batiment quand l'administration
%  qu'il abrite est devenue francaise.
%
%  LE VIOLON (81) EST « الكمان », ET C'EST L'HARMONIUM GUJARATI.
%  Cinq instruments a l'alinea 6 : violon, flute (82), violoncelle
%  (83), trombone (84), tambour (85). Trois sont transcrits. Deux ont
%  un nom d'ici, et ce sont les deux du takht, l'ensemble de la
%  musique arabe : « الكمان » et « الناي ». Or le violon n'est pas
%  plus indigene ici qu'ailleurs — il est arrive d'Europe au XIXe
%  siecle et a remplace la rababa. IL EST DEVENU ARABE COMME
%  L'HARMONIUM EST DEVENU GUJARATI : par l'usage, et en une
%  generation. Le tambour, lui, n'a jamais eu besoin de le devenir.
%
%  LA CATHEDRALE (20) ET SON CLOCHER (13) NE DEMANDENT AUCUN EFFORT A
%  CETTE COLONNE, ET C'EST RARE DANS LE RELEVE. Beyrouth, Damas et
%  Alep ont leurs cathedrales et leurs cloches dans la meme rue que
%  leurs minarets, et la langue a les deux series de mots sans que
%  l'une soit un emprunt de l'autre — كاتدرائية et جامع, ناقوس et
%  مئذنة. La colonne gujaratie avait du ecrire દેવળ pour une chose
%  qui n'est pas a tous les coins de rue.
%
%  Enfin la porte (8) : « الباب » n'est pas plus un mot d'archeologie
%  ici qu'a Ahmedabad. A Damas, Bab Touma, Bab Charqi et Bab al-Jabiya
%  sont des adresses ou l'on descend du service. Le releve avait fait
%  la meme remarque au meme renvoi dans l'autre colonne : DEUX VILLES
%  MURAILLEES, DEUX LANGUES, ET LES MEMES PORTES ENCORE DEBOUT DANS
%  LES CONVERSATIONS.
% ===================================================================

%%K t11-apar-1 apar
\VUsaut{2.0mm}
\VUcentre{13.2pt}{90}{السلسلة التالتة}
\VUsaut{3.0mm}
\VUfilet{16mm}
\VUsaut{5.6mm}
\VUtitre{15.0pt}{60}{اللوحة رقم ١١}
\VUsaut{1.4mm}
\VUpk{11.4pt}{40}{المدينة وآثارها. --- الحريق.}
\VUsaut{2.2mm}

%%K t11-tit-1 sub
\VUpk{10.2pt}{40}{الضاحية.}
\VUsaut{3.0mm}

%%K t11-c1-01-1 p
1. --- أنا ساكن بمدينة كبيرة على بعد ١٠٠ كيلومتر عن المحيط، ومع هيك
هيّي \VUgras{مرفأ} \textsuperscript{(1)} من الدرجة الأولى. والنهر اللي عم
يحدّها عرضه ٤٠٠ متر وعم يعمل خليج كتير حلو.

%%K t11-c1-02-1 p
2. --- وكل القسم من المدينة اللي عالـ\VUgras{أرصفة} \textsuperscript{(2)}
شكله بحري وصناعي تماماً، وعم يعمل \VUgras{ضاحية} \textsuperscript{(3)}
ساكنينها بس البحّارة والعمّال. وهدول عم يشتغلوا بـ\VUgras{المعامل}
\textsuperscript{(4)} وبـ\VUgras{معمل الغاز} \textsuperscript{(5)} اللي
\VUgras{مدخنته} \textsuperscript{(6)} العالية و\VUgras{خزّان الغاز} تبعه
بينشافوا من كتير بعيد.

%%K t11-c1-03-1 p
3. --- وبالقرون الوسطى، كانت المدينة محصّنة، ولهلّق منلاقي كم أثر من
أسوارها، وخصوصاً \VUgras{باب} \textsuperscript{(8)} \VUgras{القلعة}
\textsuperscript{(7)} العتيقة المحصّنة، اللي على جنبيه \VUgras{برجين}
\textsuperscript{(9)} عم يستعملوهن هلّق \VUgras{سجن}.

%%K t11-c1-04-1 p
4. --- وبكل هالـ\VUgras{حيّ}، اللي شكله كتير عتيق، في بس بنا واحد
حديث نسبياً : هوّي \VUgras{الجمرك} \textsuperscript{(10)}. وهوّي بين
الرصيف و\VUgras{الزاروب} \textsuperscript{(11)} اللي بيوصل لـ\VUgras{دار
البلدية} \textsuperscript{(12)} العتيقة. وسهل تعرف هالأثر الأخير من
\VUgras{برج الأجراس} \textsuperscript{(13)} العملاق اللي عم يطلّ على
المدينة القديمة.

%%K t11-c1-05-1 p
5. --- وعالجهة التانية من الشارع، \VUgras{واقف} بنا متل الجمرك
بالشكل : هوّي \VUgras{البورصة} \textsuperscript{(14)}. ووحدة من
واجهاته عم تطلّ على ساحة زغيرة فيها \VUgras{سراي المحافظة}
\textsuperscript{(15)}. والحقيقة إنّه مدينتنا، مع إنّها مش عاصمة، هيّي
مركز محافظة مهمّ.

%%K t11-tit-2 sub
\VUsaut{5.0mm}
\VUpk{10.2pt}{40}{أعلى قسم من المدينة.}
\VUsaut{3.4mm}

%%K t11-c1-06-1 p
6. --- والحامية، اللي عددها لا بأس فيه، ساكنة بـ\VUgras{ثكنات}
\textsuperscript{(16)} واسعة وكتير صحّية ؛ وهيّي ممتدّة لحدّ سياج
\VUgras{الحديقة العامّة} \textsuperscript{(17)}. و\VUgras{الجادّة}
\textsuperscript{(19)} اللي بتوصل لهالحديقة بتمرق وراء \VUgras{صندوق
التوفير} \textsuperscript{(18)} وبتنتهي بساحة \VUgras{الكاتدرائية}
\textsuperscript{(20)}.

%%K t11-c1-07-1 p
7. --- وكل هالآثار ممتدّة متل مدرّج على تلّة زغيرة، وين كمان
\VUgras{ثانويتنا} \textsuperscript{(21)} الواسعة.

%%K t11-c1-07-2 p
وإذا نزل الواحد بطريق شوي مايلة، بتوصل عالشارع الكبير، بيوصل
لـ\VUgras{المحكمة} \textsuperscript{(22)} اللي قدّامها ممتدّة
\VUgras{حديقة عامّة} \textsuperscript{(26)} حلوة. و\VUgras{ساحة الحديقة}
مش كتير كبيرة، بس بتعمل بنصّ المدينة واحة خضرة وبرودة. لهيك أهل
المدينة بيجوها كتير، اللي بيحبّوا يقعدوا تحت خيمة \VUgras{المقهى}
\textsuperscript{(25)}، ليسمعوا موسيقيي الجيش اللي بيعزفوا يوم الخميس
ويوم الأحد بـ\VUgras{جوسق الموسيقى} \textsuperscript{(27)} المحطوط بين
المروج وأحواض الزهر.

%%K t11-c1-08-1 p
8. --- وعلى مرج منهن واقف \VUgras{تمثال} \textsuperscript{(28)} برونز،
أثر انبنى لذكرى واحد من أهل بلدنا، اللي عمل خدمات كبيرة لمدينته اللي
انولد فيها.

%%K t11-tit-3 sub
\VUsaut{4.0mm}
\VUpk{10.2pt}{40}{وسايل النقل.}
\VUsaut{3.4mm}

%%K t11-c1-09-1 p
9. --- لهلّق مدينتنا لسّا ما تنوّرت بالضو الكهربائي، عندها بس
\VUgras{عيون الغاز} \textsuperscript{(29)}. بس \VUgras{جرايد البلد} عم
تعمل حملة ليتحسّن هالنظام تبع الإنارة، \VUgras{متل} ما
\VUgras{نظام جرّ الترامواي} انتحسّن من قبل. و\VUgras{العربايات}
\VUgras{العتيقة}، \VUgras{المجرورة} بـ\VUgras{الخيل}، انبدلت عملياً
بـ\VUgras{الترامواي الكهربائي} \textsuperscript{(31)} اللي عم ينقل
المسافرين بسرعة \VUgras{من طرف المدينة للتاني}، بسعر كتير
\VUgras{رخيص}.

%%K t11-c1-10-1 p
10. --- وجنب المحكمة في \VUgras{البنك} \textsuperscript{(32)} وين
بينحسموا السندات التجارية. و\VUgras{محصّلي البنك} \textsuperscript{(33)}
بيروحوا يقبضوا الديون عالبيوت.

%%K t11-tit-4 sub
\VUsaut{4.0mm}
\VUpk{10.2pt}{40}{الفنّ والتعليم.}
\VUsaut{3.4mm}

%%K t11-c1-11-1 p
11. --- والبنا الحلو اللي واقف عن قريب هوّي \VUgras{المتحف}. وفيه
سوا \VUgras{مكتبة} \textsuperscript{(34)} المدينة، و\VUgras{مجموعات
العلوم الطبيعية} \textsuperscript{(35)} الحلوة (متحف الطبيعة)،
و\VUgras{متحف رسم} \textsuperscript{(36)} رشيق بيعمل \VUgras{معرض} دائم
رائع.

%%K t11-c1-11-2 p
وبيفتح للناس مرّتين بالأسبوع، بس الغربا فيهن يزوروه أي يوم مقابل
بخشيش زغير لـ\VUgras{الحارس} \textsuperscript{(37)}.
و\VUgras{الرسّامين} \textsuperscript{(38)} بيجوا لهونيك ليشتغلوا.
منشوفهن \VUgras{قدّام} حاملهن، و\VUgras{لوحة الألوان} بإيدهن، عم ينسخوا
اللوحات المشهورة.

%%K t11-c1-12-1 p
12. --- وعالمرتفع، وراء المتحف، بتبلّش تبيّن \VUgras{قبّة}
\textsuperscript{(40)} \VUgras{المستشفى} \textsuperscript{(39)} وأبنية
\VUgras{دار المعلّمين} \textsuperscript{(41)} اللي كان بيتدرّب فيها
معلّمين مدارسنا البلدية.

وبساحة المسرح في \VUgras{مكتب البريد} \textsuperscript{(43)} مركّب
بـ\VUgras{بيت خاصّ} \textsuperscript{(42)}.

%%K t11-tit-5 sub
\VUsaut{5.0mm}
\VUpk{11.4pt}{40}{الحريق.}
\VUsaut{3.0mm}

%%K t11-tit-6 sub
\VUpk{10.2pt}{40}{صرخة الحريق.}
\VUsaut{3.4mm}

%%K t11-c2-01-1 p
1. --- خبر مرعب عم ينتشر بالمدينة : هلّق صار حريق بالمسرح ! وباللحظة
اللي وصلت فيها عالـ\VUgras{ساحة} \textsuperscript{(96)}، كان الجمهور من
هلّق متجمّع عالـ\VUgras{رصيف} \textsuperscript{(46)} وعالـ\VUgras{طريق}
\textsuperscript{(47)}. و\VUgras{موظّفي البريد} \textsuperscript{(44)}
و\VUgras{السعاة} \textsuperscript{(45)} عم يطلعوا من مكتب البريد.
و\VUgras{سائق ترامواي} \textsuperscript{(48)} اضطرّ يوقّف عربايته
ليخلّي \VUgras{سيّارة الإسعاف} \textsuperscript{(50)} البلدية تمرق، اللي
جايي مع \VUgras{جرّاح} \textsuperscript{(52)} و\VUgras{ممرّض}
\textsuperscript{(51)} ليعالجوا \VUgras{مجروح} \textsuperscript{(53)}
انجاب على \VUgras{نقّالة} \textsuperscript{(54)} جنب \VUgras{كشك
الجرايد} \textsuperscript{(55)}.

%%K t11-c2-02-1 p
2. --- وسريّة مشاة، كانت راجعة من التمارين، وقفت لتحفظ النظام مع
\VUgras{الدرك البلدي عالخيل} \textsuperscript{(61)}. و\VUgras{الفرسان}،
اللي خيلهن عم تتنطّط، أحسن من \VUgras{العساكر} \textsuperscript{(57)} أو
\VUgras{الدرك} \textsuperscript{(63)}، عم يرجّعوا \VUgras{الفضوليين}
\textsuperscript{(56)} لورا. وبعدين الدرك كتير مشغولين.
وواحد منهن عم يدلّ \VUgras{الشرطة} \textsuperscript{(64)} وين لازم ينقلوا
\VUgras{مصاب} \textsuperscript{(65)} تاني، إطفائي شجاع وقع من على السلّم.

%%K t11-c2-03-1 p
3. --- \VUgras{الضابط} \textsuperscript{(66)} عم يحدّد المحلّ اللي لازم
يحطّوا فيه \VUgras{مضخّة البخار} \textsuperscript{(67)} اللي جايي. ولحدّ
ما توصل هالمكنة القوية، خلّاهن يركّبوا بسرعة \VUgras{مضخّة اليد}
\textsuperscript{(68)} اللي \VUgras{خرطومها} \textsuperscript{(69)} الطويل
عم يرمي الميّ متل السيل عالنار.

وستّ إطفائيين عم يشغّلوها، وقت رفيقهن، اللي ماسك \VUgras{الفوّهة}
\textsuperscript{(70)}، عم يوجّه \VUgras{النافورة} \textsuperscript{(71)}
لقلب الحريق.

%%K t11-c2-04-1 p
4. --- وعالجهة التانية من المسرح، سندوا \VUgras{السلّم}
\textsuperscript{(72)} الكبير. وهوّي عم يخلّي الإطفائيين يقرّبوا من
اللهب اللي عم يطلع بين دوّامات \VUgras{الدخان} \textsuperscript{(75)}
الأسود.

%%K t11-tit-7 sub
\VUsaut{5.0mm}
\VUpk{10.2pt}{40}{أعمال بطولة.}
\VUsaut{3.4mm}

%%K t11-c2-05-1 p
5. --- \VUgras{الحريق} \textsuperscript{(74)} بلّش ببهو \VUgras{المسرح}
\textsuperscript{(73)}، وقت كانوا عم يتمرّنوا على \VUgras{الأوبرا} اللي
\VUgras{عرضها} الأول كان رح يصير بكرا. ولحسن الحظّ ما كان في غير كم
\VUgras{متفرّج} \textsuperscript{(76)} بالقاعة. والمساكين التجوا
عالـ\VUgras{شرفة} \textsuperscript{(77)} وين \VUgras{الإطفائية}
\textsuperscript{(86)} الشجعان جايين ينقذوهن بالسلّم وبالحبال.

%%K t11-c2-06-1 p
6. --- وتحت \VUgras{الرواق} \textsuperscript{(78)}، وين
\VUgras{الإعلان} \textsuperscript{(79)} عم يخبّر عن العرض، منشوف
\VUgras{عازفين} \textsuperscript{(81)} \VUgras{الجوقة}
\textsuperscript{(80)} عم يهربوا وهنّي شايلين آلاتهن :
\VUgras{كمان} \textsuperscript{(81)}، و\VUgras{ناي} \textsuperscript{(82)}،
و\VUgras{تشيلو} \textsuperscript{(83)}، و\VUgras{ترومبون}
\textsuperscript{(84)}، و\VUgras{طبل} \textsuperscript{(85)}، وهيك. وعم
يجتهدوا ينقذوا جزء من \VUgras{الأثاث} \textsuperscript{(87)}.

%%K t11-c2-07-1 p
7. --- و\VUgras{الممثّلين} \textsuperscript{(89)}
و\VUgras{الممثّلات} \textsuperscript{(88)}، و\VUgras{المغنّيين}
و\VUgras{المغنّيات} اضطرّوا يهربوا بـ\VUgras{ثياب المسرح}
\textsuperscript{(90)} تبعهن. والتينور عم يشرح لـ\VUgras{صحفي}
\textsuperscript{(92)} كيف صار هالشي. و\VUgras{المراسل} إجا يركض من أول
لحظة مع أصحاب السلطة : \VUgras{المحافظ} \textsuperscript{(91)}،
و\VUgras{رئيس البلدية} \textsuperscript{(93)}، و\VUgras{مفوّض الشرطة}
\textsuperscript{(94)} و\VUgras{المحقّق العدلي} \textsuperscript{(95)}،
اللي رح يحقّقوا ليحكموا بالمسؤوليات.

\VUsaut{6.0mm}
\VUfilet{22mm}
